\documentclass[10pt,a4paper]{amsart}
\usepackage[margin=19mm]{geometry}
\usepackage{amsmath,amssymb,mathtools}
\usepackage[scr=rsfs,cal=euler]{mathalfa}
\usepackage{xfrac}
\usepackage[unicode,hidelinks,bookmarks=false]{hyperref}
\newcommand{\ie}{i.e.\ }
\newcommand{\cf}{cf.\ }
\newcommand{\resp}{respectively\ }
\newcommand{\Subsection}{\subsection}
\providecommand{\nobreakdash}{\nobreak\hbox{-}}
\setlength{\emergencystretch}{2em}
\multlinegap0pt
\usepackage{bm}
\usepackage{bbm}
\usepackage{dsfont}
\usepackage{xspace}
\usepackage{cancel}
\usepackage{mathtools}
\usepackage{amsthm}
\makeatletter
\@ifundefined{thm}{\newtheorem{thm}{Theorem}}{}
\@ifundefined{rem}{\newtheorem{rem}[thm]{Remark}}{}
\makeatother
\title[Sperner's colorings of hypergraphs arising from edgewise tri]{Sperner's colorings of hypergraphs arising from edgewise triangulations}
\author{Du\v{s}ko Joji\'c, Ognjen Papaz}
\date{Source: arXiv:2506.07201v1 (CC BY 4.0)}
\begin{document}
\maketitle

\noindent\emph{Source and licence.} Sperner's colorings of hypergraphs arising from edgewise triangulations. Version arXiv:2506.07201v1 (CC BY 4.0), distributed under the Creative Commons Attribution 4.0 International licence, as recorded in the arXiv licence metadata for this exact version and shown on the abstract page https://arxiv.org/abs/2506.07201v1. Full rights notice: licenses/arxiv-2506.07201-CC-BY-4.0.txt.\par\smallskip

\noindent\textit{Editorial scope.} Counted unit: the Thm whose source label is \texttt{T:imadvadobra} in the article (source0.tex lines 1515--1518, source label \texttt{T:imadvadobra}), delivered together with the article's own complete original proof of it (source0.tex lines 1519--1595, one \texttt{proof} environment of 77 lines). No mathematics is written, repaired or re-derived: statement and proof are the article's own text line for line. Required context retained uncounted: R:blocksinpi (lines 269--279); R:Blocks (lines 1110--1124). Every definition, display and label of those lines is retained unchanged. Omitted from this package and not redistributed: 1--268, 280--1109, 1125--1514, 1596--1886. Nothing inside a retained excerpt is omitted or altered: every retained body is delivered exactly as the article's lines are written, so no omission receipt of any kind accompanies this package. Citations: the retained excerpts carry no citation command at all, so no bibliography entry of the article is transcribed and no reference is redistributed. Reference adaptation: none was needed and none was made; every reference inside the delivered unit resolves inside it, so no unresolved reference or citation is produced. Printed slips kept literal: no printed word, symbol, hypothesis or number of the counted statement or of its proof was altered. Layout and class: the article's own class is \texttt{birkjour} with packages including hyperref, color, bbm; the wrapper is the lane's pinned \texttt{amsart} editorial wrapper at 10pt on A4, which restates the theorem-like environments the retained text needs and adds the article's own macro definitions that the retained text needs (none), with extra wrapper packages (bm, bbm, dsfont, xspace, cancel, mathtools). The delivered numbering is the unit's own. Assets: no retained span contains a figure environment, a graphics-inclusion command, a PDF-page inclusion, a table or a TikZ picture; no image file is copied into this package, the package declares no asset dependency and no image file is redistributed with it. Licence: version arXiv:2506.07201v1 of this article is distributed under the Creative Commons Attribution 4.0 International licence, as recorded in the arXiv licence metadata for this exact version and shown on the abstract page https://arxiv.org/abs/2506.07201v1. A full rights notice is delivered at licenses/arxiv-2506.07201-CC-BY-4.0.txt. Characters: every retained excerpt is pure ASCII, so the delivered engine, XeLaTeX with its default Latin Modern text font, reports no missing character.
\begin{rem}\label{R:blocksinpi} A \textit{block} in a
permutation $\pi\in\mathbb{S}_{k-1}$
is a maximal subsequence of $\pi$ consisting of consecutive
integers from $[k-1]$. If $S(\pi)=\{a_1,a_2,\ldots,a_{m-1}\}$,
then
 $\pi$ has $m$ blocks
$$B_1=\{1,2,\ldots,a_1\},\ldots,
B_{i}=\{a_{i-1}+1 ,\ldots,a_i\},\ldots, B_{m}=\{a_{m-1}+1
,\ldots,k-1\}.$$ Note that the first and the last element of each
block are determined by $S(\pi)$.
\end{rem}

\begin{rem}\label{R:Blocks}
\label{P:blocksinr-invariant} Let $m<k$ be an integer relatively
prime to $k$. Assume that $k=mt+s$ for some integers $t \geq 0$
and $0<s<m$. Define a permutation $\pi \in \mathbb{S}_{k-1}$ as
follows: set $\pi_m=1$, and $\pi_{im}=i$ for all $i = 2, \ldots, k
- 1$, where the indices $im$ are taken modulo $k$. Note that $\pi$
is an $r$-invariant permutation, and in particular we have
$\pi_{k-m}=k-1$, and $\pi_s=k-t$.

Furthermore, $\pi$ has exactly $m$ almost equicardinal blocks.
Namely, $s-1$ of these blocks (those starting with
$\pi_1,\ldots,\pi_{s-1}$) are of length $t+1$, and the remain
$m-s+1$ of them is of length $t$. Since the block $B_1$ (starting
with $1$) has $t$ elements, we conclude that $\pi_{m-s}=t+1$
\end{rem}

\begin{thm}\label{T:imadvadobra} A permutation
$\pi\in \mathbb{S}_{k-1}$ has a good two-element set of colors if
and only if $\pi$ is not $r$-invariant.
\end{thm}

\begin{proof}
It follows from Remark \ref{R:Blocks} that an $r$-invariant
permutation does not admit a good $2$-element set of colors. Let
$\pi=\pi_1\pi_2\ldots \pi_{k-1}\in \mathbb{S}_{k-1}$ be a
permutation that does not admit $\{a,b\}$ as a good set of colors,
for any choice of $a$ and $b$. Assume that $\pi$ has $m$ blocks,
i.e., $S(\pi)=\{a_1,a_2,\ldots,a_{m-1}\}$. The blocks of $\pi$ are
determined by $S(\pi)$, see Remark \ref{R:blocksinpi}.
\begin{itemize}
   \item[(1)] For all $i>1$ the block $B_i$ begins before $B_1$
   ; that is $a_{i}+1$ appears before $1$ in $\pi$.
    Otherwise, if $a_{i}+1$ appears between $1$ and $a_i$, then
    $\{1,a_{i}+1\}$ is a good set for $\pi$.
    \item[(2)] If a block $B_i$ (with $i>1$) has more than one element,
     then $1$ lies
    between $a_{i}+1$ and $a_{i}+2$ (the first and second
    elements of $B_i$).
    Otherwise, $\{1,a_{i}+2\}$ is a good set for $\pi$.
\end{itemize}
In a similar way, considering $k$ as a possible color, we conclude
that
 \begin{itemize}
    \item[(3)] For all $i<m$ the block $B_i$ ends after $B_m$,
   that is $k-1$ appears before all $a_i\in
    S(\pi)$.
     \item[(4)] If a block $B_i$ (for $i<m$) has more than one element,
      then $k-1$ lies
    between $a_{i}-1$ and $a_{i}$ (the last two
    elements of $B_i$).
\end{itemize}
From $(1)-(4)$ it follows that
     \begin{itemize}
    \item[(5)] There are no two elements of a block $B_i$ appearing
    before the
    beginning of any other block $B_j$.
    \item[(6)] There are no two elements of a block $B_i$ appearing
    after the end
    of any other block $B_j$.
\end{itemize}
Since $\{a,b\}$ is not a good set for the permutation $\pi$, it
follows that $\pi$ does not contain a subsequence of the form
$$\cdots (a-1)\cdots (b-1)\cdots b \cdots a \cdots \textrm{ or }
\cdots (b-1)\cdots (a-1) \cdots a \cdots b\cdots. $$ Therefore, we
conclude the following
  \begin{itemize}
    \item[(7)] There are no two elements of a block $B_i$ that appear
    between
    two consecutive elements of another block $B_j$.
    \end{itemize}
    From $(1)-(7)$ we conclude that the blocks of $\pi$ are almost
    equicardinal.
    Therefore, if $k=tm+s$, then $s-1$ blocks of $\pi$ have $t+1$
elements, and $m-s+1$ blocks have $t$ elements. Furthermore, from
$(1)$ and $(4)$, we know that $\pi_m=1$ and $\pi_{s}=k-t$.

Since $\{a,b\}$ is not a good set for the permutation $\pi$, it
follows that $\pi$ does not contain a subsequence of the form
$$\cdots a\cdots b\cdots (b-1) \cdots (a-1)\cdots \textrm{ or}
\cdots b\cdots a\cdots (a-1) \cdots (b-1)\cdots. $$ Therefore, we
conclude the following
\begin{itemize}
    \item[(8)] If a block $B_i$ ends after $B_j$, then $B_{i+1}$ begins
    after
    $B_{j+1}$. Similarly, if $B_i$ ends before $B_j$,
    then $B_{i+1}$ begins
before
    $B_{j+1}$.
\end{itemize}

Note that the block $B_1$ ends before exactly $s-1$ blocks
(namely, all blocks of length $t+1$). From $(8)$, it follows that
the block $B_2$ begins before exactly $s-1$ blocks. Therefore, we
deduce that $\pi_{m-s}=t+1$. Continuing in the same manner, we
conclude that $\pi$ is the $r$-invariant permutation described in
Remark \ref{R:Blocks}.

\end{proof}

\end{document}
